\documentclass[11pt,reqno]{amsart}
\usepackage[T1]{fontenc}
\usepackage{lmodern}
\usepackage{amsmath,amssymb,mathtools,booktabs,array}
\usepackage{microtype}
\usepackage[margin=1.05in]{geometry}
\usepackage[hidelinks]{hyperref}
\hypersetup{pdftitle={A bound of 212 for gaps between consecutive primes},
 pdfauthor={OpenAI GPT-5.6 Sol},
 pdfsubject={Bounded gaps between primes; a computer-assisted proof},
 pdfkeywords={prime gaps, Maynard-Tao sieve, interval arithmetic}}
\usepackage{enumitem}
\setlist[enumerate]{label=\textup{(\roman*)},leftmargin=2.2em}
\setlength{\parskip}{0pt}
\setlength{\emergencystretch}{1em}
\newtheorem{theorem}{Theorem}[section]
\newtheorem{proposition}[theorem]{Proposition}
\newtheorem{lemma}[theorem]{Lemma}
\newtheorem{corollary}[theorem]{Corollary}
\theoremstyle{definition}
\newtheorem{definition}[theorem]{Definition}
\theoremstyle{remark}
\newtheorem{remark}[theorem]{Remark}
\numberwithin{equation}{section}
\newcommand{\PP}{\mathbb P}
\newcommand{\NN}{\mathbb N}
\newcommand{\RR}{\mathbb R}
\newcommand{\ind}{\mathbf 1}
\newcommand{\dd}{\,\mathrm d}
\newcommand{\Jb}{J_{\mathrm{box}}}
\newcommand{\Qset}{\mathcal Q}
\newcommand{\Tset}{\mathcal T}
\newcommand{\HH}{\mathcal H}
\newcommand{\tmax}{\tau_0}
\newcommand{\gmin}{\gamma_-}
\newcommand{\gmax}{\gamma_+}
\DeclareMathOperator{\diam}{diam}
\title[A bound of 212 for prime gaps]{A bound of 212 for gaps between consecutive primes}
\author{OpenAI GPT-5.6 Sol}
\thanks{This is a substantially AI-generated mathematical manuscript.
Siddhartha Mahajan, Research Fellow at Lossfunk, supplied the problem
and source materials, gave brief high-level directions, and arranged its
release. He did not independently derive or verify the mathematical argument.
The research and manuscript were
generated in the ChatGPT web interface by GPT-5.6 Sol using Pro reasoning
effort; Codex was used afterward for review and computational verification.
See Appendix~\ref{app:provenance} for the full provenance and release status.
The public repository is
\url{https://github.com/Siddhartha-Mahajan/prime-gap-212}. Human
correspondence: Siddhartha Mahajan,
\href{mailto:siddharthamahajan03@gmail.com}
{\texttt{siddharthamahajan03@gmail.com}}.}
\date{September 2, 2026}

\begin{document}
\begin{abstract}
We prove that every admissible $45$-tuple has infinitely many translates
containing at least two primes. In particular,
$\liminf_{n\to\infty}(p_{n+1}-p_n)\le 212$.
The distributional input is Stadlmann's equidistribution estimates for
moduli with a large smooth factor. A three-part partition lemma enlarges
the permissible support, and a sum of three product weights with separate
radial profiles supplies the sieve function. Its integrals are reduced to
finite convolutions and prefix sums. A directed-rounding certificate gives
$45J_{\mathrm{box}}/I>1.0001949699$; the box numerator is a lower bound for
the required sieve integral. The exact coefficients and a separately
implemented verifier are included in the ancillary material and at
\url{https://github.com/Siddhartha-Mahajan/prime-gap-212}.
\end{abstract}
\maketitle

\section{Introduction}
Let $p_n$ denote the $n$th prime, and put
\[
 H_1=\liminf_{n\to\infty}(p_{n+1}-p_n).
\]
The sieve of Goldston, Pintz and Y{\i}ld{\i}r{\i}m \cite{GPY} provided the
starting point for the modern theory of small prime gaps. Zhang's
theorem \cite{Zhang} established that $H_1$ is finite. Maynard's
multidimensional sieve \cite{Maynard}, together with Tao's independent
work and the subsequent Polymath project \cite{PolymathSieve}, led to
$H_1\le246$. Stadlmann \cite{Stadlmann26} obtained $H_1\le240$ by combining
the epsilon-enlargement of the sieve support with equidistribution
estimates for moduli having a large smooth factor. These estimates allow
some nonsmooth factors to remain, provided that they can be allocated to
appropriate divisor ranges.

We prove the following result.
\begin{theorem}\label{thm:main}
One has $H_1\le212$.
\end{theorem}
A finite set of integers is \emph{admissible} if it omits a residue class
modulo every prime. We obtain Theorem~\ref{thm:main} from the stronger
statement below.
\begin{theorem}\label{thm:tuples}
Let $\{h_1,\ldots,h_{45}\}$ be an admissible set of distinct integers.
There are infinitely many integers $n$ for which at least two of
$n+h_1,\ldots,n+h_{45}$ are prime.
\end{theorem}

There are two ingredients beyond the sieve and distribution estimates
of \cite{Stadlmann26}. The first is a partition argument for the
logarithms of the nonsmooth factors. The support is described by four
bounds: one for a single large coordinate, one for two, one for three,
and one for four or more. We retain the dependence between the available
divisor capacities, rather than replacing each capacity by its separate
minimum. This gives a uniform three-part partition in the most
restrictive Type~II range. A direct factorization argument retains a
positive auxiliary slack throughout this range.

The second ingredient is the choice and evaluation of the sieve
function. We use three product weights, each multiplied by a function of
the sum of the coordinates. After restricting this function to a union
of boxes, both quadratic forms can be expressed through one-dimensional
convolutions. The function on this union is a finite, exactly specified
object. In particular, the proof does not require an error estimate for
approximating an unspecified continuous optimizer, nor does it require
certifying an eigenvalue calculation.

The distribution theorems themselves are not strengthened here. We use
Lemmas~3, 4, 6, and~7 of \cite{Stadlmann26}, which extend estimates from
\cite{BakerIrving,PolymathDistribution,Stadlmann25}. The role of the new
partition lemma is to place a larger family of moduli within their
hypotheses. We use a nonnegative prime-supported weight of asymptotic
density one, so that no loss from a prime minorant enters the variational
inequality.  A direct Heath--Brown reduction at the endpoint
$\xi_2=\xi_3=2/5$ supplies the required prime-distribution estimate.

Sections~\ref{sec:sieve}--\ref{sec:distribution} establish the permissible
support. Section~\ref{sec:discrete} gives the finite quadratic forms.
Section~\ref{sec:certificate} records the certificate and proves the two
theorems. The finite certificate is specified in Appendix~\ref{app:certificate},
and the generation history and release status are recorded in
Appendix~\ref{app:provenance}.

\subsection{Notation}
All integrals are Lebesgue integrals. A list may be empty; its sum is then
zero. The entries in the lists used for factorization are real numbers.
We write $P(x)=\prod_{p<x}p$, and say that an integer is $x^\delta$-smooth
if all of its prime factors are less than $x^\delta$. All moduli to which
we apply a distribution estimate are squarefree. The parameters
$\varepsilon$, $\eta$, $\varepsilon_0$, and $\tau$ have different roles:
$\varepsilon$ enlarges the sieve support, $\eta=10^{-10}$ is the tolerance
in the convolution decomposition, $\varepsilon_0$ contracts the support
when sieve weights are constructed, and $\tau>0$ is the auxiliary
parameter in the distribution estimates.

\section{The sieve criterion}\label{sec:sieve}
Let $k\ge2$, and suppose that
\begin{equation}\label{eq:support-conditions}
 0<\varepsilon<A<\frac12-\varepsilon,\qquad
 0<\delta<B_m\le B_{m+1}\le B_m+\delta\quad(m\ge1),\quad B_1<1.
\end{equation}
For $\mathbf t=(t_1,\ldots,t_k)\in[0,\infty)^k$, define
\[
 r(\mathbf t)=\#\{i:t_i>\delta\}.
\]
We use the support
\begin{equation}\label{eq:support}
 \Tset_k=\left\{\mathbf t\in[0,\infty)^k:\ \sum_i t_i<A+\varepsilon,\quad
 \sum_{i:t_i>\delta}t_i\le B_{r(\mathbf t)}\right\},
\end{equation}
where the second condition is omitted when $r(\mathbf t)=0$.
This is the one-layer case of \cite[Definition~1]{Stadlmann26}, with
$\underline A=(-\varepsilon,A)$.

A useful consequence of \eqref{eq:support-conditions} is the following
inheritance property. If a list of $r$ numbers, all at least $\delta$,
has sum at most $B_r$, then every sublist of $j$ numbers has sum at most
$B_j$. Indeed, deleting $r-j$ entries decreases the sum by at least
$(r-j)\delta$, whereas
\begin{equation}\label{eq:inheritance}
 B_r\le B_j+(r-j)\delta.
\end{equation}
In particular, each coordinate of \eqref{eq:support} is at most $B_1$,
apart from immaterial boundary conventions.

For a symmetric real function $F\in L^2(\Tset_k)$, extended by zero, set
\begin{align}
 I(F)&=\int_{[0,\infty)^k} F(\mathbf t)^2\dd\mathbf t,
       \label{eq:I}\\
 J(F)&=\int_{\substack{\mathbf t\in[0,\infty)^{k-1}\\
                  t_1+\cdots+t_{k-1}\le A-\varepsilon}}
       \left(\int_0^\infty F(\mathbf t,u)\dd u\right)^2
       \dd\mathbf t.\label{eq:J}
\end{align}

We specify the required family of moduli. For $0<\varepsilon_0<1$, let
$\Qset(x;\varepsilon_0)$ consist of the squarefree integers admitting a
factorization
\begin{equation}\label{eq:qfactor}
 q=ee'\prod_{i=1}^{m}f_i\prod_{i=1}^{m'}f'_i
\end{equation}
such that $ee'$ is $x^\delta$-smooth, $\log_x f_i,\log_x f'_i\ge\delta$,
and
\begin{align}
 \log_x\prod_i f_i&\le(1-\varepsilon_0)B_m,&
 \log_x\prod_i f'_i&\le(1-\varepsilon_0)B_{m'},\label{eq:qcaps}\\
 \log_x\left(e\prod_i f_i\right)&\le(1-\varepsilon_0)(A-\varepsilon),&
 \log_x\left(e'\prod_i f'_i\right)&\le(1-\varepsilon_0)(A+\varepsilon).
 \label{eq:qtotals}
\end{align}
A condition on an empty product in \eqref{eq:qcaps} is omitted. Thus
\begin{equation}\label{eq:qmaximum}
 q\le x^{2A(1-\varepsilon_0)}.
\end{equation}
This is the squarefree part of the one-layer family in
\cite[Definition~2]{Stadlmann26}.

For a function $f(n;x)$ write
\[
 E_f(x;q,a)=\sum_{\substack{x\le n\le2x\\n\equiv a\pmod q}}f(n;x)
 -\frac1{\phi(q)}\sum_{\substack{x\le n\le2x\\(n,q)=1}}f(n;x).
\]
For positive integers $n$ define the nonnegative prime-supported weight
\begin{equation}\label{eq:rho}
 \rho_x(n)=\frac{\log n}{\log(3x)}\ind_\PP(n)
             \ind_{x/2\le n\le3x}.
\end{equation}
Thus $0\leq\rho_x\leq\ind_\PP$ on the range relevant to the sieve, and
the prime number theorem gives
\[
 \sum_{x\leq n\leq2x}\rho_x(n)=(1+o(1))\frac{x}{\log x}.
\]
The distribution statement needed below is
\begin{equation}\label{eq:prime-distribution}
 \sum_{q\in\Qset(x;\varepsilon_0)}|E_{\rho_x}(x;q,a)|
 \ll_{C,\varepsilon_0}\frac{x}{(\log x)^C}
 \qquad\text{for every }C>0,
\end{equation}
with the residue convention $(a,P(x))=1$ of
\cite[Definition~3]{Stadlmann26}.

\begin{proposition}[One-layer sieve criterion]\label{prop:sieve}
Suppose \eqref{eq:support-conditions} and
\eqref{eq:prime-distribution} hold. If a symmetric real
$F\in L^2(\Tset_k)$ satisfies
\begin{equation}\label{eq:sieve-threshold}
 kJ(F)>I(F)>0,
\end{equation}
then every admissible $k$-tuple has infinitely many translates
containing at least two primes.
\end{proposition}
\begin{proof}
Apply \cite[Proposition~1 and its proof]{Stadlmann26} with one support
layer and $\rho=\rho_x$. The minorant parameters are $c_1=c_2=0$.
The roughness condition is automatic because $\rho_x$ is supported on
primes, and its mean was recorded above. In the one-layer case the
integrals in that proposition are exactly \eqref{eq:I} and
\eqref{eq:J}. The sieve argument applies to any fixed admissible tuple.

The stated $L^2$ formulation includes nonsmooth functions. One can also
see the relevant continuity directly: integration in the last coordinate
has squared operator norm at most $B_1$, by Cauchy--Schwarz. Hence $J$ is
continuous on $L^2(\Tset_k)$. The contraction and smooth approximation in
\cite[Lemma~1]{Stadlmann26} therefore preserve the strict inequality
\eqref{eq:sieve-threshold}.
\end{proof}

\section{Partition lemmas}\label{sec:partition}
The next two lemmas concern two lists of positive real numbers. A
$j$-subsum means the sum of any $j$ distinct entries from one list, not
necessarily consecutive entries.

\begin{lemma}\label{lem:two-bins}
Let $0<b_1\le b_2\le b_3$, and suppose that every $j$-subsum of either
of two lists is at most $b_j$, for $1\le j\le3$. Let $a,c>0$ satisfy
\begin{equation}\label{eq:large-items}
 2b_1\le a,\qquad \frac52b_2\le a+c,\qquad b_3\le3c.
\end{equation}
The sum of all entries exceeding $c$ is at most $a$. Moreover, if the
combined total is $M$ and $2(M-a)\le c$, the entries can be partitioned
into parts of sums at most $a$ and $c$.
\end{lemma}
\begin{proof}
There are at most two entries exceeding $c$ in each list, since three
such entries would have sum greater than $b_3$. If each list has at most
one, their combined sum is at most $2b_1\le a$. Otherwise one list has
two entries exceeding $c$, so $c<b_2/2$. The sum of all entries exceeding
$c$ is then at most $2b_2$, and
\[
 a\ge\frac52b_2-c>2b_2.
\]
This proves the first assertion.

If $M\le a$, there is nothing to prove. Otherwise put $R=M-a>0$.
An entry in $[R,c]$ can be placed in the second part on its own. If no
such entry exists, every entry at most $c$ is less than $R$. The first
assertion implies that these smaller entries have total at least $R$.
Accumulate them until their sum first reaches $R$. The resulting sum is
less than $2R\le c$, and the sum of the remaining entries is at most $a$.
\end{proof}

\begin{lemma}[Three-part partition]\label{lem:three-bins}
Let $0<\delta<b_1\le b_2\le b_3\le B$. Suppose that all entries in
two lists are at least $\delta$, each list has total at most $B$, and
every $j$-subsum of either list is at most $b_j$, for $1\le j\le3$.
Let $a,c,d>0$, and put $c'=\max(c,d)$, $d'=\min(c,d)$.
Suppose that
\begin{align}
 2b_1&\le a,& \frac52b_2&\le a+c',& b_3&\le3c',\label{eq:three-one}\\
 2B&\le a+d',& 4B&\le2a+c'+2\delta.\label{eq:three-two}
\end{align}
Then the entries can be partitioned into three parts of sums at most
$a,c,d$, respectively.
\end{lemma}
\begin{proof}
Let $M$ be the combined total. We may assume that $M>a$.
Suppose first that there is an entry $y\le d'$. Place $y$ in the part of
capacity $d'$. Deleting this entry preserves all the subsum conditions.
The remaining total $M'=M-y$ satisfies
\[
 2(M'-a)\le 2(2B-\delta-a)\le c'.
\]
Lemma~\ref{lem:two-bins} partitions the remaining entries into the parts
of capacities $a$ and $c'$.

If there is no entry at most $d'$, then every entry exceeds $d'$.
By \eqref{eq:three-two}, $R=M-a\le2B-a\le d'$. By the first assertion of
Lemma~\ref{lem:two-bins}, the entries exceeding $c'$ have total at most
$a$. Since $M>a$, some entry is at most $c'$. This entry is greater than
$d'\ge R$, so placing it in the part of capacity $c'$ leaves total at
most $a$. The part of capacity $d'$ is empty.
\end{proof}

For completeness, we record the elementary fact used to turn a partition
of the large factors into a factorization of the modulus.
\begin{lemma}[Filling an interval]\label{lem:fill}
Let $z_1,\ldots,z_s\in(0,\delta)$, and let $Y\ge0$.
If $U-L\ge\delta$, $Y\le U$, and $Y+\sum_i z_i\ge L$, some sublist has
sum $Z$ with $Y+Z\in[L,U]$. If instead $Y<U$ and
$Y+\sum_i z_i>L$, one can obtain $Y+Z\in(L,U)$ whenever $U-L>\delta$.
\end{lemma}
\begin{proof}
If the initial value is already in the desired interval, use the empty
sublist. Otherwise add entries until the lower endpoint is first reached,
or strictly exceeded in the open-interval version. The last increment is
less than $\delta$.
\end{proof}

\section{The critical divisor factorization}\label{sec:factorization}
We now fix the parameters used in the proof. All entries in the following
display are exact:
\begin{equation}\label{eq:parameters}
\begin{gathered}
 k=45,\qquad
 A=\frac{16862}{65536},\quad
 \varepsilon=\frac{557}{65536},\quad
 \delta=\frac{1019}{65536},\\
 b_1=\frac{10171}{65536},\quad
 b_2=\frac{10378}{65536},\quad
 b_3=\frac{11321}{65536},\quad
 B=\frac{12082}{65536},\\
 B_1=b_1,\qquad B_2=b_2,\qquad B_3=b_3,\qquad B_m=B\quad(m\ge4),\\
 \xi_1=\frac{19}{50},\qquad \xi_2=\xi_3=\frac25,\qquad
 \eta=10^{-10},\quad \tmax=10^{-8},\quad \zeta=10^{-6}.
\end{gathered}
\end{equation}
These parameters satisfy \eqref{eq:support-conditions}. Put
\begin{equation}\label{eq:gamma-parameters}
\begin{gathered}
 \omega=A-\frac14=\frac{239}{32768},\qquad
 G=\frac13+8\omega+\frac73\delta=\frac{28047}{65536},\\
 \gmin=\xi_2-\eta,\qquad \gmax=G+3\eta.
\end{gathered}
\end{equation}
We will take $\tau>0$ sufficiently small, always with $\tau\le\tmax$.
Define
\begin{align}
 \alpha&=\gmin-2\delta-8\omega-109\tmax,\label{eq:alpha}\\
 S&=\frac12-2\delta-10\omega-116\tmax,\label{eq:S}\\
 C_*&=\max\left(\frac12-\gmax-2\omega-7\tmax,
                    \frac45\left(\frac12-\gmax\right)\right).
                    \label{eq:Cstar}
\end{align}
Direct rational comparison gives
\begin{align}
 2b_1&<\alpha,& \frac52b_2&<S,& b_3&<3C_*,\label{eq:uniform-one}\\
 2B&<S,& 2B&<\gmin-2\delta-4\tmax,&
 4B&<S+\alpha+2\delta.\label{eq:uniform-two}
\end{align}
Quantitative margins are listed in Appendix~\ref{app:margins}.

\begin{proposition}\label{prop:critical-factor}
Let $q\in\Qset(x;\varepsilon_0)$ satisfy $q\ge x^{1/2-2\tau}$.
For every $\gamma\in[\gmin,\gmax]$, there are divisors $r\mid q$,
$u\mid q/r$, and $d_1\mid r$ such that, with $\delta_* =\delta+\zeta$,
\begin{align}
 x^{\gamma-3\tau-\delta_*}&<r<x^{\gamma-3\tau},\label{eq:r-window}\\
 x^{1-\gamma-6\tau-\delta_*}q^{-1}&<u<
                 x^{1-\gamma-6\tau}q^{-1},\label{eq:u-window}\\
 r^2x^{2-\gamma-52\tau-\delta_*}q^{-4}&<d_1<
                 r^2x^{2-\gamma-52\tau}q^{-4}.\label{eq:d-window}
\end{align}
\end{proposition}
\begin{proof}
Write $q=x^{1/2+2\omega_0}$ exactly. By \eqref{eq:qmaximum},
$-\tau\le\omega_0\le\omega$. Apply the inheritance property
\eqref{eq:inheritance} to the two lists of logarithms in
\eqref{eq:qfactor}. Each list has total at most $B$, and its $j$-subsums
are bounded by $b_j$ for $1\le j\le3$.

Set
\begin{equation}\label{eq:capacities}
 a=\gamma-2\delta-8\omega_0-109\tau,\qquad
 c=\frac12-\gamma-2\omega_0-7\tau,\qquad
 d=8\omega_0+105\tau.
\end{equation}
These numbers are positive; in particular $d\ge97\tau$.
Writing $c'=\max(c,d)$ and $d'=\min(c,d)$, we have
\begin{gather}
 a\ge\alpha,\qquad a+c\ge S,\qquad
 a+d=\gamma-2\delta-4\tau\ge\gmin-2\delta-4\tmax,
 \label{eq:capacity-sums}\\
 c'\ge c\ge\frac12-\gmax-2\omega-7\tmax,\qquad
 c'\ge\frac{4c+d}{5}
      =\frac45\left(\frac12-\gamma\right)+\frac{77}{5}\tau.
 \label{eq:capacity-max}
\end{gather}
Thus $c'\ge C_*$. Moreover
\[
 a+d'=\min(a+c,a+d),\qquad
 2a+c'\ge a+(a+c)\ge\alpha+S.
\]
Equations~\eqref{eq:uniform-one}--\eqref{eq:uniform-two} verify all
hypotheses of Lemma~\ref{lem:three-bins}. We obtain three parts of the
large logarithms with respective sums
\begin{equation}\label{eq:Yparts}
 Y_1\le a,\qquad Y_2\le c,\qquad Y_4\le d.
\end{equation}
The indices $1,2,4$ identify their roles in the factorization; no fourth
nonempty part is used.

Write the prime-factor logarithms of the smooth part $ee'$ as
$z_1,\ldots,z_s\in(0,\delta)$, and put
\begin{align*}
 \ell_1&=\gamma-4\tau-\delta,&u_1&=\gamma-4\tau,\\
 \ell_2&=\frac12-\gamma-2\omega_0-7\tau-\delta,&u_2&=c,\\
 \ell_3&=-\gamma-8\omega_0-101\tau-\delta,&
 u_3&=-\gamma-8\omega_0-101\tau.
\end{align*}
Each interval has length $\delta$, and
\begin{equation}\label{eq:window-identities}
 2\ell_1+u_3=a,\qquad
 \ell_1-2u_1-\ell_3=d,\qquad
 a+d=\gamma-2\delta-4\tau.
\end{equation}

First construct $u$ from the large factors in the second part and some
smooth primes. The initial logarithm is at most $u_2$, and the total
available logarithm exceeds $\ell_2$, since
\[
 \frac12+2\omega_0-(a+d)-\ell_2
   =4\omega_0+3\delta+11\tau\ge3\delta+7\tau>0.
\]
Lemma~\ref{lem:fill} gives $\log_xu\in[\ell_2,u_2]$.
Next construct $r$ from the large factors in the first and third parts
and some of the remaining smooth primes. Their initial logarithm is
at most $a+d=u_1-2\delta$, and the total available logarithm is at least
\[
 \frac12+2\omega_0-u_2=\gamma+4\omega_0+7\tau
 \ge\gamma+3\tau>\ell_1.
\]
Thus $v=\log_xr\in[\ell_1,u_1]$ can be obtained.

Finally, use the large factors in the first part and smooth primes
already selected for $r$. The initial logarithm is at most
$a=2\ell_1+u_3\le2v+u_3$, while the total available logarithm is at least
\[
 v-Y_4\ge\ell_1-d=2u_1+\ell_3\ge2v+\ell_3.
\]
A further application of Lemma~\ref{lem:fill} gives a divisor $d_1\mid r$
with
\[
 2v+\ell_3\le\log_xd_1\le2v+u_3.
\]
All selections use disjoint prime factors where required, since $q$ is
squarefree.

The closed intervals used for $r$ and $u$ lie strictly inside
\eqref{eq:r-window} and \eqref{eq:u-window}; their lower-endpoint slack is
$\zeta-\tau>0$ and their upper-endpoint slack is $\tau>0$. For the last
divisor, write
\[
 T=2\log_xr+2-\gamma-4\log_xq.
\]
The constructed interval is
\[
 [T-101\tau-\delta,T-101\tau].
\]
It lies strictly inside \eqref{eq:d-window}, with lower and upper slacks
$\zeta-49\tau>0$ and $49\tau>0$. This proves the proposition.
\end{proof}

\section{Verification of the distribution hypotheses}\label{sec:distribution}
We use the coefficient-sequence and Siegel--Walfisz conventions of
\cite[Definitions~6--9]{Stadlmann26}.  We shall only need the Type~II
and Type~III members of the convolution family in Definition~9 of that
paper.  Its Type~II sequences have two Siegel--Walfisz factors, one at
scale $x^\gamma$ with
$\xi_2-\eta\le\gamma\le1-\xi_2+\eta$. The Type~III scales satisfy the
conditions in that definition. We retain all of its coefficient bounds
and uniformity conventions.

The following reduction derives the required prime-supported estimate from
the Type~II and Type~III distribution estimates.

\begin{lemma}[Endpoint Heath--Brown reduction]\label{lem:HB-reduction}
Let $\mathcal R(x)$ be a family of squarefree moduli satisfying
\[
 x^{1/2-2\tau}\le q\le x^{1/2+2\omega}\qquad(q\in\mathcal R(x)).
\]
Suppose that, for every primitive residue class $a$, with arbitrary
logarithmic saving and uniformly in the coefficient sequences, the
discrepancy sum over $\mathcal R(x)$ is bounded by $x(\log x)^{-C}$ for
\begin{enumerate}
\item every Type~II convolution in $\HH(\xi_1,2/5,2/5)$, and
\item every Type~III convolution in $\HH(\xi_1,2/5,2/5)$.
\end{enumerate}
If $\omega<1/20$ and $0<\tau\le10^{-8}$, then the same conclusion
holds for $\rho_x$.
\end{lemma}
\begin{proof}
We adapt the proof of \cite[Lemma~2.7 and Section~3]
{PolymathDistribution}, retaining only the alternatives that are needed
at the endpoint $2/5$.  Put
\[
 \sigma=\frac1{10}+\frac\eta4.
\]
Thus $\sigma>1/10$ and, for our parameters, $\sigma>2\omega$.
Apply the Heath--Brown identity with $K=10$, for which $1/K<2\sigma$,
followed by the finer-than-dyadic smooth decomposition
used in \cite[Section~3]{PolymathDistribution}.  Up to
$O((\log x)^{O(1)})$ terms, whose number is absorbed by the arbitrary
logarithmic saving, this reduces the von Mangoldt function on $[x,2x]$
to convolutions
\[
 \alpha_1\star\cdots\star\alpha_s
\]
at scales $N_i=x^{t_i+o(1)}$, where $t_i\ge0$ and
$\sum_i t_i=1$.  By \cite[Lemma~3.4]{PolymathDistribution}, every
subconvolution at a scale $x^c$ with fixed $c>0$ has the
Siegel--Walfisz property, and every individual factor at scale at least
$x^{2\sigma}$ is smooth.  Choose the smooth decomposition sufficiently
fine that all $o(1)$ endpoint errors below are smaller than $\eta/4$.

Apply \cite[Lemma~3.1]{PolymathDistribution} to the $t_i$.  In its
Type~I/II alternative there is a partition into two groups, the smaller
of which has exponent $\gamma$ satisfying
\[
 \frac12-\sigma+o(1)<\gamma\le\frac12.
\]
Both grouped convolutions are Siegel--Walfisz, and
\[
 \frac12-\sigma=\frac25-\frac\eta4>\frac25-\eta.
\]
Consequently this is a Type~II convolution from the family in the
statement.

In the Type~III alternative there are three distinct smooth factors
with exponents $t_i,t_j,t_k$ satisfying
\[
 2\sigma\le t_i,t_j,t_k\le\frac12-\sigma,
 \qquad t_i+t_j,t_i+t_k,t_j+t_k\ge\frac12+\sigma.
\]
The identities
\[
 2\sigma=\frac15+\frac\eta2,qquad
 \frac12-\sigma=\frac25-\frac\eta4,qquad
 \frac12+\sigma=\frac35+\frac\eta4
\]
place this convolution strictly inside the Type~III ranges for
$\xi_3=2/5$, even after the endpoint errors.  It is therefore covered
by the second hypothesis.

It remains to treat the Type~0 alternative. Write the convolution as
\[
 F=\alpha\star\beta,\qquad
 \alpha(m)=\psi(m/N),\qquad
 N\ge x^{1/2+\sigma+o(1)},\qquad M\asymp x/N,
\]
where $\beta$ is the complementary coefficient sequence at scale $M$ and
$\psi$ has the smoothness bounds furnished by Lemma~3.4 of that paper.
For primitive $a\pmod q$, terms with $(\ell,q)>1$ vanish from both sides
of the discrepancy, and hence
\begin{multline*}
 E_F(x;q,a)=\sum_{(\ell,q)=1}\beta(\ell)\left(
 \sum_{\substack{x/\ell\le m\le2x/\ell\\
                  m\equiv a\bar\ell\pmod q}}\alpha(m)\right.\\
 \left.{}-\frac1{\phi(q)}
 \sum_{\substack{x/\ell\le m\le2x/\ell\\(m,q)=1}}\alpha(m)
 \right).
\end{multline*}
For fixed $\ell$, the function
\[
 g_\ell(t)=\psi(t/N)\ind_{[x/\ell,2x/\ell]}(t)
\]
has supremum and total variation bounded by $(\log x)^{O(1)}$.
Euler summation in a residue class, and M\"obius inversion for the
coprimality condition, therefore bound the expression in parentheses by
\[
 \tau(q)(\log x)^{O(1)}=x^{o(1)}.
\]
Since $\sum_\ell|\beta(\ell)|\ll M(\log x)^{O(1)}$ and
$\#\mathcal R(x)\le x^{1/2+2\omega}$, we obtain
\[
 \sum_{q\in\mathcal R(x)}|E_F(x;q,a)|
 \ll x^{1/2+2\omega}M(\log x)^{O(1)}
 \ll x^{1-\sigma+2\omega+o(1)}.
\]
This is $O_C(x(\log x)^{-C})$ for every $C$, because
$\sigma>2\omega$. We have proved the required distribution estimate for
$\Lambda$ with the cutoff included.

Finally remove prime powers.  On primes,
$\Lambda(n)=\log n$, while the difference is supported on $p^j$ with
$j\ge2$.  Uniformly for primitive $a$, squarefreeness gives at most
$2^{\#\{p:p\mid q\}}$ roots to $z^2\equiv a\pmod q$.  Hence the total progression
contribution of prime squares, summed over
$x^{1/2-2\tau}\le q\le x^{1/2+2\omega}$, is
\[
 \ll x^{1/2+2\omega+2\tau+o(1)}.
\]
For $j\ge3$, the crude bound of $O(x^{1/3})$ possible bases for each
modulus gives
\[
 \ll x^{5/6+2\omega+o(1)}.
\]
The averaged main terms contribute only $x^{1/2+o(1)}$, using
$\sum_{q\le Q}\phi(q)^{-1}\ll\log Q$.  All three bounds are
$O_C(x(\log x)^{-C})$ for every $C$.  Dividing by $\log(3x)$ and using
\eqref{eq:rho} proves the assertion for $\rho_x$.
\end{proof}

\begin{proposition}\label{prop:distribution}
The parameters \eqref{eq:parameters} satisfy
\eqref{eq:prime-distribution}.
\end{proposition}
\begin{proof}
Fix $\varepsilon_0>0$ and a required logarithmic saving. Choose $\tau>0$
sufficiently small for the distribution estimates below, with
$\tau\le\tmax$. The subfamily $q<x^{1/2-2\tau}$ is covered directly,
for $\rho_x$, by the Bombieri--Vinogradov theorem for $\Lambda$, removal
of prime powers, and division by $\log(3x)$. We therefore
work on
\[
 \Qset^+=\{q\in\Qset(x;\varepsilon_0):q\ge x^{1/2-2\tau}\}.
\]
The logarithms of all large factors have combined sum $M\le2B$.
The smooth prime logarithms are less than $\delta$, and the total
logarithm of $q$ is at least $1/2-2\tau$.

We verify the divisor conditions of
\cite[Lemmas~3, 4, 6, and~7]{Stadlmann26}. All inequalities used for this purpose
are strict. The numerical comparisons in this proof are rational
comparisons, not floating-point tests; the exact checks are recorded in
Appendix~\ref{app:certificate}.

\subsection*{Type IIa and Type IIb}
Normalize the auxiliary scales so that $MN=x$; this changes $M$ by
only a bounded factor and preserves the coefficient-sequence and
Siegel--Walfisz properties. Interchange the two factors when necessary,
so that the smaller scale has exponent at most $1/2$, as in
\cite[Section~4.3, Case~II]{Stadlmann26}. Put
\begin{equation}\label{eq:K}
 K=\frac25+\frac{24}{5}\omega+\frac75\delta
   =\frac{149677}{327680}.
\end{equation}
For $\gamma\ge K+2\eta$, set
\[
 \Delta_a(\gamma)=\frac{5\gamma-2-24\omega}{7}-\eta.
\]
One has $\Delta_a(\gamma)-\delta\ge3\eta/7>0$, and
\begin{align*}
 24\omega+7\Delta_a(\gamma)-5\gamma&=-2-7\eta<-2,\\
 8\omega+3\Delta_a(\gamma)-\gamma
   &=\frac{8\gamma-6-16\omega}{7}-3\eta<0.
\end{align*}
These are the scalar conditions of \cite[Lemma~3]{Stadlmann26}.
The divisor interval has logarithmic endpoints
$\gamma-3\tau-\Delta_a(\gamma)$ and $\gamma-3\tau$.
Again include all large factors and fill with smooth primes. The upper
endpoint exceeds $2B$, while
\[
 \frac12-2\tau-\bigl(\gamma-3\tau-\Delta_a(\gamma)\bigr)
 \ge\frac1{14}-\frac{24}{7}\omega-\eta+\tau>0.
\]

For $G+3\eta<\gamma<K+2\eta$, put
\[
 \Delta_b(\gamma)=\frac{3\gamma-1-24\omega}{7}-\eta.
\]
Then $\Delta_b(\gamma)-\delta>2\eta/7$, and
\begin{align*}
 24\omega+7\Delta_b(\gamma)-3\gamma&=-1-7\eta<-1,\\
 8\omega+3\Delta_b(\gamma)-\gamma
   &=\frac{2\gamma-3-16\omega}{7}-3\eta<0.
\end{align*}
We apply \cite[Lemma~4]{Stadlmann26}. Choose $u$ using only smooth primes,
with logarithm in
\[
 \left(\frac12-\gamma-2\omega-6\tau-\Delta_b(\gamma),\
       \frac12-\gamma-2\omega-6\tau\right).
\]
The upper endpoint is positive, and the smooth mass exceeds even that
upper endpoint: it is enough to check
\begin{equation}\label{eq:IIb-smooth}
 G+2\omega-2B-2\tmax>0.
\end{equation}
Next choose $r$, coprime to $u$, containing all large factors, with
logarithm in
\[
 (\gamma-3\tau-\Delta_b(\gamma),\ \gamma-3\tau).
\]
Its initial logarithm is below the upper endpoint, by
$2B<G+3\eta-3\tmax$. The remaining logarithm exceeds the lower endpoint,
since
\[
 \log_x(q/u)>\gamma+2\omega+4\tau
             >\gamma-3\tau-\Delta_b(\gamma).
\]
Lemma~\ref{lem:fill} therefore produces both divisors.

\subsection*{Type IIc}
It remains to consider $\gmin\le\gamma\le\gmax$.
With $\delta_* =\delta+\zeta$, exact comparison gives
\begin{align}
 1-8\omega-4\delta_*-2\gmax&>0.02352,\label{eq:IIc-scalar-a}\\
 \gmin-32\omega-10\delta_*&>0.01110,\label{eq:IIc-scalar-b}\\
 4\gmin-1-48\omega-16\delta_*&>0.00110.\label{eq:IIc-scalar-c}
\end{align}
Thus all three scalar hypotheses of
\cite[Lemma~6]{Stadlmann26} hold throughout the interval. The required
divisors are exactly those in
Proposition~\ref{prop:critical-factor}. Notice that its argument includes
$-\tau\le\omega_0<0$ as well as $\omega_0\ge0$.

\subsection*{Type III}
Use the scale parameter $\gamma_3=\xi_3+\eta$ and put
\begin{equation}\label{eq:III-delta}
 \Delta_3=\frac12-\frac72\omega-\frac98\gamma_3-\eta.
\end{equation}
Then $\Delta_3>\delta$ and
\[
 28\omega+9\gamma_3+8\Delta_3=4-8\eta<4.
\]
The Type~III scales in $\HH(\xi_1,\xi_2,\xi_3)$ satisfy the scale
conditions of \cite[Lemma~7]{Stadlmann26} with $\gamma_3$.
Set
\begin{equation}\label{eq:III-window}
 U_3=\frac13+\frac43\Delta_3-\frac43\omega,\qquad L_3=U_3-\Delta_3.
\end{equation}
That lemma requires a divisor with logarithm in $(L_3,U_3)$.
If $M<U_3$, place all large factors in the divisor. If $M\ge U_3$,
exclude any one large factor. Its logarithm lies in $[\delta,b_1]$, so
the logarithm of the remaining large factors is at most $2B-\delta<U_3$.
In either case the available total logarithm exceeds $L_3$, since
\begin{equation}\label{eq:III-cap}
 U_3-(2B-\delta)>0.00307,\qquad
 \frac12-2\tmax-b_1-L_3>0.01303.
\end{equation}
Fill with smooth primes and apply Lemma~\ref{lem:fill}.

\subsection*{Passage to primes}
We have obtained the required distribution estimate on $\Qset^+$ for
every Type~II and Type~III convolution required in
Lemma~\ref{lem:HB-reduction}. The constants and auxiliary parameters can
be chosen uniformly on the indicated ranges: the scalar inequalities
have positive margins, and the cutoff ranges have the fixed tolerance
$\eta$. The bound \eqref{eq:qmaximum} places every modulus within
$x^{1/2+2\omega}$. Since $\omega=239/32768<1/20$,
Lemma~\ref{lem:HB-reduction} proves \eqref{eq:prime-distribution} on
$\Qset^+$.
Combining it with Bombieri--Vinogradov on the smaller moduli completes
the proof.
\end{proof}

\section{A finite sieve function}\label{sec:discrete}
We next give an exact discretization for evaluating \eqref{eq:I} and a
lower bound for \eqref{eq:J}. The formulas in this section apply to any
parameters satisfying \eqref{eq:support-conditions} that are integral
multiples of a mesh $d>0$, with the index ranges below nonempty.
For definiteness, assume $1\le D<C_1\le N+1$ and $0\le M_0\le N$.

Write
\begin{equation}\label{eq:integer-parameters}
\begin{gathered}
 D=\delta/d,\quad C_r=B_r/d,\quad
 T=(A+\varepsilon)/d,\quad L=(A-\varepsilon)/d,\\
 N=T-k,\qquad M_0=L-(k-1).
\end{gathered}
\end{equation}
For $a\ge0$ let $Q_a=[ad,(a+1)d)$. A box
$Q_{a_1}\times\cdots\times Q_{a_k}$ is retained if
\begin{equation}\label{eq:allowed-box}
 s=\sum_i a_i\le N,\qquad z+r\le C_r,
 \quad r=\#\{i:a_i\ge D\},\quad z=\sum_{a_i\ge D}a_i,
\end{equation}
where the second condition is omitted for $r=0$.
Every retained box lies in $\Tset_k$ up to a null set. The use of upper
endpoints in \eqref{eq:allowed-box} is important: a retained box is not
merely a box whose centre lies in the support.

Fix positive arrays $g_\ell(a)$ and real arrays $h_\ell(s)$, with
$0\le\ell<R$ and $0\le a,s\le N$. On a retained box define
\begin{equation}\label{eq:finite-F}
 F(t_1,\ldots,t_k)=
 \sum_{\ell=0}^{R-1}h_\ell(s)\prod_{i=1}^k g_\ell(a_i),
\end{equation}
and set $F=0$ elsewhere. This function is symmetric, bounded, and
square-integrable. No sign condition is imposed on the radial arrays.

\subsection{Convolution coefficients}
For $0\le\ell,m<R$, put
\[
 q_{\ell m}(a)=d\,g_\ell(a)g_m(a).
\]
Let $q^{\mathrm s}_{\ell m}$ be its restriction to $0\le a<D$, extended
by zero, and let $q^{\mathrm l}_{\ell m}$ be its restriction to
$D\le a<C_1$. The latter truncation is legitimate by
\eqref{eq:inheritance}. Define
\begin{align}
 P^{\ell m}_j&=(q^{\mathrm s}_{\ell m})^{*j},&
 P^{\ell m}_0(s)&=\ind_{s=0},\label{eq:small-powers}\\
 Q^{\ell m}_0(z)&=\ind_{z=0},&
 Q^{\ell m}_r(z)&=\ind_{z\le C_r-r}
             (Q^{\ell m}_{r-1}*q^{\mathrm l}_{\ell m})(z).
             \label{eq:large-powers}
\end{align}
Here $*$ denotes convolution of sequences indexed by the nonnegative
integers. In every calculation coefficients beyond the largest needed
index are discarded.

The truncation after each convolution in \eqref{eq:large-powers} loses
no term that could contribute at a later stage. Indeed, for $r\ge2$,
if $r$ large indices have sum $z\le C_r-r$, deleting an index at least $D$ leaves sum
at most
\[
 C_r-r-D\le C_{r-1}-r<C_{r-1}-(r-1).
\]
This uses $C_r-C_{r-1}\le D$. Iterating the argument reduces to $r=1$,
where the cap is $C_1-1$. It also shows directly that the individual
large indices are less than $C_1$.

\begin{proposition}\label{prop:finite-norm}
For the function \eqref{eq:finite-F},
\begin{equation}\label{eq:finite-I}
 I(F)=\sum_{\ell,m=0}^{R-1}\sum_{s=0}^{N}h_\ell(s)h_m(s)
       \sum_{r=0}^{k}\binom{k}{r}
       (P^{\ell m}_{k-r}*Q^{\ell m}_r)(s).
\end{equation}
Only values $r\le\min(k,\lfloor B/\delta\rfloor)$ can contribute when
$B_r\le B$.
\end{proposition}
\begin{proof}
Expand the square in \eqref{eq:finite-F}. Select the $r$ large
coordinates in $\binom{k}{r}$ ways. Their index sum is recorded by
$Q_r^{\ell m}$, and the sum of the remaining indices by
$P_{k-r}^{\ell m}$. Each factor $q_{\ell m}$ includes one factor $d$,
so the product already includes the volume $d^k$ of the box. Summing
over all retained boxes proves the formula.
\end{proof}

\subsection{The box numerator}
In the $(k-1)$ common coordinates of \eqref{eq:J}, retain only boxes
whose upper-endpoint sum is at most $A-\varepsilon$. Their index sum is
therefore at most $M_0$. Let $\Jb(F)$ be the integral in \eqref{eq:J}
over this union of common-coordinate boxes, with the inner fiber
integral left unchanged. Since the integrand is a square,
\begin{equation}\label{eq:box-lower}
 0\le\Jb(F)\le J(F).
\end{equation}

For a common-coordinate box with $s\le M_0$, large-coordinate count $r$,
and large index sum $z$, first impose the inherited cap $z+r\le C_r$
when $r>0$. A box failing this cap has an empty fiber. For a box satisfying
it, the allowed last-coordinate indices are exactly
$0\le a\le u(s,z,r)$, where
\begin{equation}\label{eq:fiber-cutoff}
 u(s,z,r)=\min\bigl(N-s,\ \max(D-1,C_{r+1}-r-1-z)\bigr).
\end{equation}
This follows by considering separately a small last coordinate and a
large last coordinate. Define the fiber prefixes
\begin{equation}\label{eq:prefix}
 H_\ell(s,u)=\sum_{a=0}^{u}g_\ell(a)h_\ell(s+a).
\end{equation}

\begin{proposition}\label{prop:finite-J}
For $F$ in \eqref{eq:finite-F},
\begin{multline}\label{eq:finite-J}
 \Jb(F)=d^2\sum_{\ell,m=0}^{R-1}\sum_{s=0}^{M_0}\sum_{r=0}^{k-1}
 \binom{k-1}{r}\sum_{z\ge0}
 Q^{\ell m}_r(z)P^{\ell m}_{k-1-r}(s-z)\\
 {}
 \cdot H_\ell\bigl(s,u(s,z,r)\bigr)
        H_m\bigl(s,u(s,z,r)\bigr).
\end{multline}
Terms with indices outside the arrays are zero.
\end{proposition}
\begin{proof}
The inner integral on a fixed common-coordinate box is
\[
 d\sum_{\ell=0}^{R-1}
   \left(\prod_{i=1}^{k-1}g_\ell(a_i)\right)
   H_\ell(s,u(s,z,r)).
\]
Square this expression and group the common-coordinate boxes by $s,r,z$.
Their product weights and volumes give
$\binom{k-1}{r}Q_r^{\ell m}(z)P_{k-1-r}^{\ell m}(s-z)$.
The two remaining factors of $d$ come from the two fiber integrals.
\end{proof}

\subsection{The coefficient arrays}
For \eqref{eq:parameters}, choose
\begin{equation}\label{eq:mesh}
 d=\frac1{65536},\qquad R=3.
\end{equation}
The integer data in \eqref{eq:integer-parameters} are
\begin{equation}\label{eq:integer-data}
\begin{gathered}
 D=1019,\quad T=17419,\quad L=16305,\quad N=17374,\quad M_0=16261,\\
 C_1=10171,\quad C_2=10378,\quad C_3=11321,\quad C_r=12082\ (r\ge4).
\end{gathered}
\end{equation}
At most eleven large coordinates occur. The six arrays
$g_0,g_1,g_2,h_0,h_1,h_2$, each of length $17375$, are specified in
Appendix~\ref{app:certificate}. Every entry is an exact dyadic rational;
together with \eqref{eq:allowed-box} and \eqref{eq:finite-F}, these arrays
define $F$ exactly.

The arrays were found by a variational search in the family
\eqref{eq:finite-F}. The coordinate profiles were initialized from
\[
 \left(\frac{c}{c+(a+\tfrac12)d}\right)^{p_\ell},\qquad
 c=0.0014,\quad(p_0,p_1,p_2)=(0.90,0.96,1.02),
\]
with an independent normalization for each profile. The radial arrays
were optimized on a coarser mesh and interpolated to \eqref{eq:mesh}.
This describes the search only. The definition used in the proof is the
stored dyadic array, not the displayed real-power formula or the output
of an eigenvalue solver.

\section{The certificate and the prime-gap bound}\label{sec:certificate}
\begin{proposition}[Finite certificate]\label{prop:certificate}
The function specified by \eqref{eq:finite-F},
\eqref{eq:integer-data}, and the coefficient arrays in
Appendix~\ref{app:certificate} has $I(F)>0$
and satisfies
\begin{equation}\label{eq:certified-interval}
 \frac{6037239562351545}{6036062711329456}
 \ \le\ \frac{45\Jb(F)}{I(F)}
 \ \le\ \frac{72446874766843545}{72432752528865376}.
\end{equation}
In particular,
\begin{equation}\label{eq:certified-decimal}
 1.0001949699<\frac{45\Jb(F)}{I(F)}<1.0001949704.
\end{equation}
\end{proposition}
\begin{proof}
Evaluate \eqref{eq:finite-I} and \eqref{eq:finite-J} by directed
rounding, with the input entries interpreted as exact dyadic rationals.
The interval computation performs the convolutions
\eqref{eq:small-powers}--\eqref{eq:large-powers} directly and evaluates the
prefix sums \eqref{eq:prefix}. It returns lower and upper bounds for each
pair $(\ell,m)$ with $\ell\le m$. Off-diagonal contributions are included
with multiplicity two.

The polynomial coefficients are nonnegative. Their lower and upper
bounds are therefore obtained by downward and upward rounding,
respectively, at every operation. The factors involving $h_\ell$ may
have either sign; these are enclosed by interval multiplication and
addition. The resulting hexadecimal endpoints are converted to exact
rational numbers before the six pair contributions are combined.
In particular, the exact aggregate enclosures are
\begin{align*}
 \frac{2263523516527043}{2251799813685248}
 &\le I(F)\le\frac{1131761758374273}{1125899906842624},\\
 \frac{402482637490103}{18014398509481984}
 &\le\Jb(F)\le\frac{1609930550374301}{72057594037927936}.
\end{align*}
The norm has a strictly positive lower bound. Dividing the lower
numerator by the upper denominator, and conversely for the upper bound,
gives \eqref{eq:certified-interval}.
The arithmetic procedure and complete endpoint data are specified in
Appendix~\ref{app:certificate}.
\end{proof}

\begin{proof}[Proof of Theorem~\ref{thm:tuples}]
Proposition~\ref{prop:distribution} verifies the distribution hypothesis
of Proposition~\ref{prop:sieve}. The function $F$ is symmetric and belongs
to $L^2(\Tset_{45})$. By
\eqref{eq:box-lower} and Proposition~\ref{prop:certificate},
\[
 45J(F)\ge45\Jb(F)>I(F)>0.
\]
Proposition~\ref{prop:sieve} proves the assertion.
\end{proof}

\begin{proof}[Proof of Theorem~\ref{thm:main}]
Consider the following set $\mathcal A$:
\begin{equation}\label{eq:tuple}
\begin{split}
\{&0,2,12,14,24,26,30,32,36,50,54,56,60,66,72,74,80,84,\\
  &92,96,102,110,114,116,122,126,134,140,144,150,156,162,\\
  &164,170,176,180,182,186,192,194,200,204,206,210,212\}.
\end{split}
\end{equation}
\begin{samepage}
It has $45$ distinct elements and diameter $212$. The table below gives
one residue class omitted by $\mathcal A$ for each prime at most $45$.
\begin{center}
\small
\begin{tabular}{@{}c*{14}{r}@{}}
\toprule
$p$&2&3&5&7&11&13&17&19&23&29&31&37&41&43\\
Omitted&1&1&3&6&9&3&1&6&17&13&28&31&8&25\\
\bottomrule
\end{tabular}
\end{center}
\end{samepage}
For $p>45$, a $45$-element set cannot occupy all residue classes.
Thus $\mathcal A$ is admissible. By Theorem~\ref{thm:tuples}, infinitely
many of its translates contain two primes. Between any such pair there
is a consecutive-prime gap at most $212$. These pairs occur arbitrarily
far out, proving $H_1\le212$.
\end{proof}

\section*{Acknowledgments}
Siddhartha thanks Paras Chopra, founder
of \href{https://lossfunk.com}{Lossfunk}, for introducing him to the general
research direction of using
AI systems to attack mathematical problems. He also thanks Lossfunk, a
Bangalore-based research lab, for providing the research workspace,
environment, and GPT access that enabled this project. Siddhartha also
thanks Julia Stadlmann, whose work
\cite{Stadlmann26} this manuscript heavily builds upon.

\appendix
\section{Finite certificate}\label{app:certificate}\label{app:margins}
The arrays in \eqref{eq:finite-F} are specified by
\path{ancillary/witness.hex}. Its first line contains
$(k,R,N+1,D,d^{-1},T,L,r_{\max})$, and its second line contains
$(C_0,\ldots,C_{r_{\max}+1})$, where $k=45$, $C_0=0$, and $r_{\max}=11$;
the remaining parameter values are given in
\eqref{eq:mesh}--\eqref{eq:integer-data}.
For each $a=0,\ldots,N$, the next line lists
$g_0(a),g_1(a),g_2(a),h_0(a),h_1(a),h_2(a)$ as exact hexadecimal
dyadic rationals. The SHA--256 digest, with the two lines concatenated, is
\begin{center}
\small\ttfamily
4b82ce9794faeb80d1766dddc2884d6a\\
19734490baf15a4200bb15bdfd5af718
\end{center}

The verifier \path{ancillary/certify.cpp} evaluates
\eqref{eq:finite-I} and \eqref{eq:finite-J} using IEEE~754 binary64
arithmetic with gradual underflow and correctly rounded addition,
multiplication, and division. The input dyadics and all required binomial
coefficients are represented exactly. Each operation is rounded outwards;
signed products use the minimum and maximum of the four endpoint
products. Induction therefore gives an enclosure of every intermediate
quantity. The six pair enclosures are combined with exact rational
arithmetic, giving \eqref{eq:certified-interval}. The computation assumes
that the implementation honors the rounding modes and compiler options
\begin{verbatim}
-O3 -std=c++17 -fopenmp -frounding-math
-ffp-contract=off -fno-fast-math
\end{verbatim}

With Python~3 and a C++17 compiler supporting this arithmetic model and
OpenMP, the certificate is reproduced by
\begin{verbatim}
python3 ancillary/run_certificate.py --threads 4
\end{verbatim}
This checks the witness digest, the rational support inequalities, and
tuple admissibility, then compiles the verifier and evaluates the finite
sums. The recorded pair endpoints are in
\path{ancillary/reference_4threads.json}.

The six margins in \eqref{eq:uniform-one}--\eqref{eq:uniform-two}
exceed, in their displayed order,
\[
 (0.0001576,\ 0.0000751,\ 0.0001434,\ 0.02725,\ 0.0001891,\ 0.0001869).
\]
These terminating decimals are strict rational lower bounds. All
parameter comparisons in Section~\ref{sec:distribution} are checked by
\path{ancillary/verify_support.py}; their exact margins are recorded in
\path{ancillary/support_output.json}.

\section{Research provenance and release status}
\label{app:provenance}

\subsection{Generation and human involvement}
This manuscript was generated in the ChatGPT web interface by
\href{https://developers.openai.com/api/docs/models/gpt-5.6-sol}
{OpenAI GPT-5.6 Sol} (model identifier \texttt{gpt-5.6-sol}), using Pro
reasoning effort. Siddhartha Mahajan, Research Fellow at Lossfunk, supplied
the bounded-prime-gaps problem, directed GPT-5.6 Sol to the relevant source
materials, and gave brief, high-level instructions,
including the instruction to seek the strongest reduction it could support.
GPT-5.6 Sol generated the principal mathematical argument, parameter search,
manuscript text, certificate data, and verification software. Codex was used
afterward to inspect the release package, rerun its computational checks, and
assist with verification and preparation of the final text.

There has been no human verification of this work to date.

The model is named as the author of this manuscript to reflect how
the work was produced.

\subsection{Public archive}
The manuscript source, ancillary code, exact witness, reference outputs,
checksums, and a human-readable provenance record are available at
\begin{center}
\url{https://github.com/Siddhartha-Mahajan/prime-gap-212}.
\end{center}

\begin{thebibliography}{9}
\bibitem{BakerIrving}
R. C. Baker and A. J. Irving,
\emph{Bounded intervals containing many primes},
Math. Z. \textbf{286} (2017), 821--841.

\bibitem{GPY}
D. A. Goldston, J. Pintz, and C. Y. Y\i ld\i r\i m,
\emph{Primes in tuples I},
Ann. of Math. (2) \textbf{170} (2009), 819--862.

\bibitem{Maynard}
J. Maynard,
\emph{Small gaps between primes},
Ann. of Math. (2) \textbf{181} (2015), 383--413.

\bibitem{PolymathDistribution}
D. H. J. Polymath,
\emph{New equidistribution estimates of Zhang type},
Algebra Number Theory \textbf{8} (2014), no.~9, 2067--2199.

\bibitem{PolymathSieve}
D. H. J. Polymath,
\emph{Variants of the Selberg sieve, and bounded intervals containing many primes},
Res. Math. Sci. \textbf{1} (2014), 1--83.

\bibitem{Stadlmann25}
J. Stadlmann,
\emph{On primes in arithmetic progressions and bounded gaps between many primes},
Adv. Math. \textbf{468} (2025), 110190.

\bibitem{Stadlmann26}
J. Stadlmann,
\emph{Bounded gaps between primes},
arXiv:2608.31126v1 [math.NT], 2026.

\bibitem{Zhang}
Y. Zhang,
\emph{Bounded gaps between primes},
Ann. of Math. (2) \textbf{179} (2014), no.~3, 1121--1174.
\end{thebibliography}
\end{document}
